%%% DODATKI DO SZABLONU WIT
%%% Plik settings.tex zostaje bez zmian; wszystko, co dodane, jest tutaj.

%%% 1. TABELE
\RequirePackage{array}
\RequirePackage{booktabs}
\RequirePackage{tabularx}
\RequirePackage{multirow}
\RequirePackage{longtable}
\newcolumntype{L}{>{\raggedright\arraybackslash}X}
\newcolumntype{C}{>{\centering\arraybackslash}X}

% Zalecenia edytorskie Wydziału: tytuł tabeli centralnie nad tabelą, podpis rysunku
% centralnie pod rysunkiem, tabela 10 pkt. Przy klasie 12 pkt to \footnotesize: każda tabela
% ma je wpisane sama, bo LaTeX zeruje rozmiar czcionki na początku wstawki.
\captionsetup{justification=centering}
\captionsetup[table]{position=top, skip=6pt}
\AtBeginEnvironment{longtable}{\footnotesize}

% Domyślne położenie: klocki nie podają go same, więc w katalogu mogą stać w miejscu.
\floatplacement{table}{htbp}
\floatplacement{figure}{htbp}

%%% 2. POLSKIE NAZWY DLA \cref (settings.tex definiuje tylko rysunki, tabele, równania, listingi, algorytmy)
\crefname{chapter}{rozdz.}{rozdz.}
\Crefname{chapter}{Rozdział}{Rozdziały}
\crefname{section}{podrozdz.}{podrozdz.}
\Crefname{section}{Podrozdział}{Podrozdziały}
\crefname{subsection}{podrozdz.}{podrozdz.}
\Crefname{subsection}{Podrozdział}{Podrozdziały}
\crefname{appendix}{zał.}{zał.}
\Crefname{appendix}{Załącznik}{Załączniki}
\crefname{subfigure}{Rys.}{Rys.}
\Crefname{subfigure}{Rysunek}{Rysunki}
\Crefname{listing}{Listing}{Listingi}
\crefname{lstlisting}{listing}{listingi}
\Crefname{lstlisting}{Listing}{Listingi}

%%% 3. STYL KODU C#
%%% literate zastępuje ustawienie z settings.tex, więc są tu też polskie znaki.
\lstdefinestyle{csharp}{
  language=[Sharp]C,
  morekeywords={var, get, set, yield, where, nameof, unsafe, stackalloc, ref},
  % 10 pkt jak w tabelach; długi wiersz łamie się na spacji, a ciąg dalszy ma strzałkę.
  basicstyle=\ttfamily\footnotesize,
  breakatwhitespace=true,
  postbreak=\mbox{\textcolor{gray}{$\hookrightarrow$}\space},
  literate={ą}{{\k{a}}}1 {ć}{{\'c}}1 {ę}{{\k{e}}}1 {ł}{{\l{}}}1 {ń}{{\'n}}1 {ó}{{\'o}}1 {ś}{{\'s}}1 {ż}{{\.z}}1 {ź}{{\'z}}1
           {Ą}{{\k{A}}}1 {Ć}{{\'C}}1 {Ę}{{\k{E}}}1 {Ł}{{\L{}}}1 {Ń}{{\'N}}1 {Ó}{{\'O}}1 {Ś}{{\'S}}1 {Ż}{{\.Z}}1 {Ź}{{\'Z}}1
           {×}{{$\times$}}1 {–}{{\textendash}}1 {−}{{$-$}}1 {·}{{$\cdot$}}1 {°}{{$^\circ$}}1 {…}{{$\ldots$}}1
           {→}{{$\rightarrow$}}1 {↔}{{$\leftrightarrow$}}1 {≈}{{$\approx$}}1 {κ}{{$\kappa$}}1 {Σ}{{$\Sigma$}}1
           {„}{{\quotedblbase}}1 {”}{{\textquotedblright}}1 {•}{{$\bullet$}}1 {■}{{$\blacksquare$}}1 {✓}{{$\checkmark$}}1
}

%%% 4. PRZEŁĄCZNIK SZKICU I MIEJSCA NA TREŚĆ
%%% \szkictrue albo \szkicfalse ustawia się w main.tex.
\newif\ifszkic
\definecolor{szkicramka}{RGB}{176,112,0}
\definecolor{szkictlo}{RGB}{255,246,224}
\newcounter{szkicwskazowki}
\newcounter{szkicpola}
\newcounter{szkicobrazy}
\newcounter{szkiclistingi}

\newcommand{\szkicRamka}[1]{%
  \par\smallskip\noindent
  \fcolorbox{szkicramka}{szkictlo}{%
    \parbox{\dimexpr\linewidth-2\fboxsep-2\fboxrule\relax}{%
      \raggedright\footnotesize\sffamily #1}}%
  \par\smallskip}

% Wskazówka w ramce; po \szkicfalse znika, a \cite w niej przestaje się liczyć do bibliografii.
\newcommand{\CoTuMaByc}[1]{%
  \stepcounter{szkicwskazowki}%
  \ifszkic\szkicRamka{\textbf{Co tu ma być:} #1}\fi}

% Pole do wpisania w tekście, tabeli albo pseudokodzie; po \szkicfalse jest puste.
\newcommand{\Wpisz}[1]{%
  \stepcounter{szkicpola}%
  \ifszkic\textcolor{szkicramka}{[#1]}\fi}

% Pusta komórka tabeli na wynik.
\newcommand{\wyn}{\Wpisz{?}}

% Nazwa kolumny CSV albo pliku z podkreśleniami.
\newcommand{\kol}[1]{\texttt{\detokenize{#1}}}

% Rozmiar siatki: E × N × piętra.
\newcommand{\siatka}[3]{#1\,$\times$\,#2\,$\times$\,#3}

% Obraz z folderu figures, bez rozszerzenia. Gdy pliku nie ma, stoi ramka zastępcza.
% [szerokość]{plik}{co ma przedstawiać}
\newcommand{\MiejsceNaObraz}[3][0.85\linewidth]{%
  \IfFileExists{figures/#2.pdf}{\includegraphics[width=#1]{figures/#2.pdf}}{%
  \IfFileExists{figures/#2.png}{\includegraphics[width=#1]{figures/#2.png}}{%
  \IfFileExists{figures/#2.jpg}{\includegraphics[width=#1]{figures/#2.jpg}}{%
    \stepcounter{szkicobrazy}%
    \fbox{\parbox[c][\dimexpr#1*9/16\relax][c]{\dimexpr#1-2\fboxsep-2\fboxrule\relax}{%
      \centering
      \ifszkic
        \footnotesize\sffamily\textbf{figures/\detokenize{#2}}\par\smallskip #3%
      \fi}}%
  }}}}

% Cały rysunek: [szerokość]{plik}{podpis}{etykieta}{co ma przedstawiać}
\newcommand{\MiejsceNaRysunek}[5][0.85\linewidth]{%
  \begin{figure}
    \centering
    \MiejsceNaObraz[#1]{#2}{#5}
    \caption{#3}
    \label{#4}
  \end{figure}}

% Listing z folderu listingi, bez rozszerzenia .cs: {plik}{podpis}{etykieta}{skąd wziąć kod}
\newcommand{\MiejsceNaListing}[4]{%
  \IfFileExists{listingi/#1.cs}{%
    \lstinputlisting[style=csharp, caption={#2}, label={#3}]{listingi/#1.cs}%
  }{%
    \stepcounter{szkiclistingi}%
    \ifszkic\szkicRamka{\textbf{Listing:} zapisz fragment jako \kol{listingi/#1.cs}. #4}\fi
    \lstinputlisting[style=csharp, caption={#2}, label={#3}]{listingi/brak.cs}%
  }}

% Po kompilacji: ile miejsc zostało do uzupełnienia (ostrzeżenie w dzienniku kompilacji).
\AtEndDocument{%
  \ifnum\numexpr\value{szkicwskazowki}+\value{szkicpola}+\value{szkicobrazy}+\value{szkiclistingi}\relax>0
    \PackageWarningNoLine{szablon pracy}{Do uzupelnienia: wskazowki \arabic{szkicwskazowki}, pola \arabic{szkicpola}, rysunki \arabic{szkicobrazy}, listingi \arabic{szkiclistingi}}%
  \fi}

%%% 5. ZALECENIA EDYTORSKIE WYDZIAŁU
%%% Szablon WIT ma większe nagłówki i interlinię 1,15. Zalecenia edytorskie: tytuł rozdziału 14 pkt,
%%% podrozdziału 13 pkt, niższego stopnia 12 pkt, interlinia pojedyncza. Włączenie: jedna linia w main.tex.
\newcommand{\WlaczZaleceniaEdytorskie}{%
  \setstretch{1}%
  \titleformat{\chapter}[display]{\bfseries\fontsize{14}{17}\selectfont}{\chaptertitlename\ \thechapter}{0.5ex}{}%
  \titleformat{\section}{\bfseries\fontsize{13}{16}\selectfont}{\thesection.}{1em}{}%
  \titleformat{\subsection}{\bfseries\fontsize{12}{15}\selectfont}{\thesubsection.}{1em}{}%
}
